\documentclass[11pt,letterpaper]{article}

\usepackage[margin=1in]{geometry}
\usepackage{titlesec}
\usepackage{enumitem}
\usepackage{hyperref}
\usepackage{booktabs}
\usepackage{longtable}
\usepackage{xcolor}
\usepackage{fancyhdr}
\usepackage{parskip}
\usepackage{tcolorbox}

\pagestyle{fancy}
\fancyhf{}
\rhead{\small EGA Access Playbook}
\lhead{\small Singlet AI}
\cfoot{\thepage}

\hypersetup{
    colorlinks=true,
    linkcolor=blue!60!black,
    urlcolor=blue!60!black,
}

\titleformat{\section}{\large\bfseries}{\thesection.}{0.5em}{}
\titleformat{\subsection}{\normalsize\bfseries}{\thesubsection.}{0.5em}{}

\newtcolorbox{actionbox}[1]{
    colback=blue!5!white,
    colframe=blue!50!black,
    title={\textbf{Action: } #1},
    fonttitle=\small,
}

\newtcolorbox{notebox}{
    colback=yellow!10!white,
    colframe=yellow!60!black,
    title={\small\textbf{Note}},
}

\title{%
    \vspace{-1.5cm}
    \textbf{EGA Controlled-Access Data}\\[0.3em]
    \large Step-by-Step Application Playbook\\[0.3em]
    \normalsize For Single-Cell Transcriptomic Reference Construction
}
\author{Singlet AI}
\date{March 2026}

\begin{document}
\maketitle
\thispagestyle{fancy}

% ══════════════════════════════════════════════════════════════════════════
\section{Overview}
% ══════════════════════════════════════════════════════════════════════════

The European Genome-phenome Archive (EGA) hosts \textbf{879 single-cell
studies} across 509 scRNA-seq, 231 10x Chromium, 58 spatial, 39 scATAC,
38 snRNA-seq, and other assay types. Unlike dbGaP's centralized
Authorized Access, EGA uses a \textbf{per-study Data Access Committee
(DAC)} model where each study's PI or institutional committee reviews
access requests independently.

This playbook provides a \textbf{prioritized, step-by-step process}
for efficiently applying to the highest-value DACs.

\subsection{Key Constraint}

We will only pursue EGA datasets for which the DAC permits
\textbf{download to university on-premises HPC infrastructure}.
Datasets requiring cloud-only processing (e.g., some UK Biobank
datasets) are excluded from this phase.

% ══════════════════════════════════════════════════════════════════════════
\section{Prerequisites}
% ══════════════════════════════════════════════════════════════════════════

Before submitting any EGA applications:

\begin{enumerate}[nosep]
    \item \textbf{Create an EGA account:}
          \url{https://ega-archive.org/register/}

    \item \textbf{Obtain institutional ethics approval:}
          Many DACs require an ethics committee opinion or IRB
          determination letter. Prepare a protocol synopsis describing
          the NMF reference construction project.

    \item \textbf{Prepare a standard Data Transfer Agreement (DTA):}
          Some European DACs require a DTA or Data Processing Agreement
          between your institution and theirs. Have your institutional
          contracts office prepare a template.

    \item \textbf{Install the EGA Download Client:}
          \begin{verbatim}
          pip install pyega3
          \end{verbatim}
          Or the newer EGA Download Client v2:
          \begin{verbatim}
          pip install ega-download-client
          \end{verbatim}

    \item \textbf{GDPR readiness:} Confirm your institution has
          Standard Contractual Clauses (SCCs) for EU--US data transfer,
          or can execute them on request.
\end{enumerate}

% ══════════════════════════════════════════════════════════════════════════
\section{Prioritized DAC Application Schedule}
% ══════════════════════════════════════════════════════════════════════════

DACs are prioritized by the number of single-cell studies they cover.
Apply to \textbf{Tier 1 first} (highest yield), then proceed to
Tier 2 and Tier 3 as Tier 1 responses arrive.

\subsection{Tier 1: Major Institutional DACs (Week 1--2)}

\begin{longtable}{@{} l l r p{4.5cm} @{}}
\toprule
\textbf{DAC} & \textbf{Institution} & \textbf{$\sim$Studies} & \textbf{Focus} \\
\midrule
\endhead
EGAC00001000205 & Wellcome Sanger Institute & 120 &
    Human Cell Atlas, cancer, immune, developmental \\
EGAC00001000654 & EMBL-EBI & 40 &
    Various European SC studies \\
EGAC00001000039 & Karolinska Institutet & 35 &
    Developmental, immune, spatial \\
EGAC00001001093 & Broad Institute & 25 &
    Cancer, immune --- may overlap with dbGaP \\
\bottomrule
\end{longtable}

\begin{notebox}
The Wellcome Sanger DAC alone covers $\sim$120 SC studies. Getting
this single approval unlocks the largest batch of EGA SC data.
Prioritize this application above all others.
\end{notebox}

\subsection{Tier 2: Specialized DACs (Week 2--4)}

\begin{longtable}{@{} l l r p{4.5cm} @{}}
\toprule
\textbf{DAC} & \textbf{Institution} & \textbf{$\sim$Studies} & \textbf{Focus} \\
\midrule
\endhead
EGAC00001000688 & Max Planck Society & 20 &
    Developmental biology, aging \\
EGAC00001001077 & University of Cambridge & 18 &
    Cancer immunogenomics, immune \\
EGAC00001000203 & Helmholtz Munich & 15 &
    Metabolic disease, lung atlas \\
EGAC00001000789 & Netherlands Cancer Inst. & 12 &
    Cancer multi-omics \\
\bottomrule
\end{longtable}

\subsection{Tier 3: Remaining DACs (Week 4+, rolling)}

The remaining $\sim$594 studies are distributed across $\sim$100+
smaller DACs (often individual PIs). For these:
\begin{itemize}[nosep]
    \item Batch by disease category or tissue to allow template reuse.
    \item Contact DAC emails listed on each study's EGA page.
    \item Expect highly variable response times (1 week to 6+ months).
    \item Some PIs may never respond --- deprioritize after two
          follow-ups.
\end{itemize}

% ══════════════════════════════════════════════════════════════════════════
\section{Step-by-Step Application Process}
% ══════════════════════════════════════════════════════════════════════════

\subsection{Step 1: Identify the DAC for Each Study}

\begin{actionbox}{Find the DAC}
For each target study (e.g., EGAS00001002072):
\begin{enumerate}[nosep]
    \item Visit \texttt{https://ega-archive.org/studies/EGAS00001002072}
    \item Note the dataset IDs (EGAD...) listed under the study
    \item Click each dataset to find the \textbf{Data Access Committee}
          field, which lists the DAC ID (EGAC...) and contact info
    \item Record the DAC ID, contact email, and any institution-specific
          requirements
\end{enumerate}
\end{actionbox}

\subsection{Step 2: Prepare the Application}

\begin{actionbox}{Customize the template}
Use the \texttt{ega\_dac\_application\_template.tex} file:
\begin{enumerate}[nosep]
    \item Replace \texttt{[DAC NAME]} with the specific DAC
    \item Replace \texttt{[DAC INSTITUTION]} with their institution
    \item Add the specific EGAS/EGAD IDs to the datasets section
    \item Fill in your institutional details (university, IRB, etc.)
    \item Compile to PDF: \texttt{pdflatex ega\_dac\_application\_template.tex}
\end{enumerate}
\end{actionbox}

\subsection{Step 3: Submit via EGA Portal}

\begin{actionbox}{Submit the request}
\begin{enumerate}[nosep]
    \item Log in to \url{https://ega-archive.org}
    \item Navigate to the dataset page (EGAD...)
    \item Click ``Request Access''
    \item Fill in the online form with information from the PDF
    \item Upload the PDF application as supporting documentation
    \item Submit
\end{enumerate}
\end{actionbox}

\begin{notebox}
Some DACs have their own application forms hosted externally (e.g.,
the Sanger Institute uses a custom portal). Check the DAC page for
institution-specific instructions before submitting.
\end{notebox}

\subsection{Step 4: Follow Up}

\begin{itemize}[nosep]
    \item If no response in 2 weeks: send a polite email to the DAC
          contact
    \item If no response in 4 weeks: send a second follow-up noting
          the original request date
    \item If no response in 8 weeks: consider the DAC non-responsive
          and deprioritize
    \item For institutional DACs (Sanger, EMBL, etc.): typical response
          is 2--4 weeks
    \item For PI-managed DACs: highly variable, often longer
\end{itemize}

\subsection{Step 5: Download Approved Datasets}

\begin{actionbox}{Download to university HPC}
Once approved, you will receive EGA download credentials:
\begin{verbatim}
# Configure pyega3
export EGA_USERNAME="your_ega_username"
export EGA_PASSWORD="your_ega_password"

# Download specific dataset
pyega3 -cf credentials.json fetch EGAD00001004149 \
    --output-dir /scratch/ega_data/

# Or use the newer client
ega-download-client download --dataset EGAD00001004149 \
    --output /scratch/ega_data/
\end{verbatim}
\end{actionbox}

\subsection{Step 6: Confirm Compliance}

After downloading:
\begin{itemize}[nosep]
    \item Verify data integrity (checksums)
    \item Confirm data is on encrypted storage
    \item Log the download in your data inventory
    \item Begin reprocessing pipeline
    \item Plan secure destruction after processing
\end{itemize}

% ══════════════════════════════════════════════════════════════════════════
\section{Paired WGS + Single-Cell Studies in EGA}
% ══════════════════════════════════════════════════════════════════════════

We identified \textbf{57 EGA studies} with paired WGS/WES and
single-cell transcriptome data. These are \textbf{top priority} for
genome--transcriptome linkage:

{\small
\begin{longtable}{@{} l p{9cm} @{}}
\toprule
\textbf{EGA ID} & \textbf{Study} \\
\midrule
\endhead
EGAS00001002072 & SC T cells from hepatocellular carcinoma patients \\
EGAS00001003387 & SC RNA-seq + WGS from different cells \\
EGAS00001006807 & Full-length SC RNA-seq with isoform + genomic alterations \\
EGAS00001002436 & Subclonal evolution of ER+ breast cancers (WGS + scRNA) \\
EGAS00001006851 & Multiregion WES + Smart-Seq3 scRNA of breast tumors \\
EGAS00001006052 & Tumor cell mutations + antigen presentation (multi-omic) \\
EGAS00001005472 & Glioma epigenetic encoding + transcriptional states \\
EGAS00001006374 & Pathogenic variants + SC transcription in disease \\
EGAS00001007078 & Clonally resolved SC multi-omics \\
\multicolumn{2}{c}{\textit{+ 48 additional studies (see catalog)}} \\
\bottomrule
\end{longtable}
}

% ══════════════════════════════════════════════════════════════════════════
\section{Protocol and Disease Distribution in EGA}
% ══════════════════════════════════════════════════════════════════════════

\begin{minipage}[t]{0.48\textwidth}
\textbf{Protocol Breakdown:}
\begin{tabular}{@{} l r @{}}
\toprule
\textbf{Protocol} & \textbf{Studies} \\
\midrule
scRNA-seq & 509 \\
10x Chromium & 231 \\
Spatial & 58 \\
scATAC-seq & 39 \\
snRNA-seq & 38 \\
Smart-seq & 31 \\
CITE-seq & 19 \\
Multiome & 10 \\
snATAC-seq & 8 \\
Drop-seq & 1 \\
\bottomrule
\end{tabular}
\end{minipage}
\hfill
\begin{minipage}[t]{0.48\textwidth}
\textbf{Disease Breakdown:}
\begin{tabular}{@{} l r @{}}
\toprule
\textbf{Category} & \textbf{Studies} \\
\midrule
Cancer & 422 \\
Neurological & 183 \\
Developmental & 157 \\
Normal/Atlas & 148 \\
Infectious & 81 \\
Pulmonary & 63 \\
Metabolic & 30 \\
Kidney & 21 \\
Autoimmune & 18 \\
Cardiovascular & 14 \\
\bottomrule
\end{tabular}
\end{minipage}

\bigskip
\noindent \textbf{Overlap with dbGaP/GEO:} Only 0.2\% of EGA studies
mention GEO and 0.1\% mention dbGaP, confirming that EGA represents a
$\sim$100\% unique data layer.

% ══════════════════════════════════════════════════════════════════════════
\section{Timeline and Tracking}
% ══════════════════════════════════════════════════════════════════════════

\begin{tabular}{@{} l l l @{}}
\toprule
\textbf{Phase} & \textbf{Action} & \textbf{Target Date} \\
\midrule
Week 1--2 & Submit Tier 1 (Sanger, EMBL, Karolinska, Broad) & $T_0$ + 2 wks \\
Week 2--4 & Submit Tier 2 (MPI, Cambridge, Helmholtz, NKI) & $T_0$ + 4 wks \\
Week 4+ & Rolling Tier 3 applications & $T_0$ + 4--12 wks \\
Ongoing & Track responses, follow up, download approved data & Continuous \\
\bottomrule
\end{tabular}

\bigskip
\noindent Maintain a shared spreadsheet tracking: DAC ID, submission date,
dataset IDs, status (submitted / approved / denied / non-responsive),
follow-up dates, download date, processing status.

\end{document}
